While \RLM{}s show strong performance on tasks beyond the context window limitations of existing LMs at reasonable inference costs, the optimal mechanism for implementing \RLM{}s remains under-explored. We focused on synchronous sub-calls inside of a Python REPL environment, but we note that alternative strategies involving asynchronous sub-calls and sandboxed REPLs can potentially significantly reduce the runtime and inference cost of \RLM{}s. Furthermore, we chose to use a max recursion depth of one (i.e. sub-calls are LMs); while we found strong performance on existing long-context benchmarks, we believe that future work should investigate deeper layers of recursion. 

Lastly, we focused our experiments on evaluating \RLM{}s using \textit{existing} frontier models. Explicitly training models to be used as \RLM{}s (e.g. as root or sub-LMs) could provide additional performance improvements -- as we found in \S\ref{sec4.4-qualitative}, current models are inefficient decision makers over their context. We hypothesize that \RLM{} trajectories can be viewed as a form of reasoning~\citep{openai2024openaio1card,deepseekai2025deepseekr1incentivizingreasoningcapability}, which can be trained by bootstrapping existing frontier models~\citep{zelikman2022starbootstrappingreasoningreasoning, zelikman2024quietstarlanguagemodelsteach}. %
